\section{Discussion}
\label{sec:discussion}
\para{{\bf interpretation.}} The central hypothesis was that explicit sentence-level ``thinking'' would increase care and reliability. Our evidence does not support that claim for this setup. The clearest effect is a \decrease in judged carefulness on \truthfulqa, while truthfulness, math accuracy, and moderation outcomes do not improve.

\para{{\bf why might this happen?}} We observe three plausible mechanisms. First, rigid tag insertion constrains response flow and can reduce natural explanatory structure. Second, visible tags may consume output budget and shorten useful content; in \truthfulqa, mean word count drops from 40.08 (\direct) to 27.14 (\thinkbetween). Third, explicit thought text may be stylistic rather than functional if the model does not use it for meaningful verification.

\para{{\bf limitations.}} This study uses one generation model family (\genmodel), one temperature setting, and 50 sampled items per task. \truthfulqa scoring relies on an LLM judge, which introduces evaluator-model dependence. The forced three-sentence style may itself influence quality independent of the thought-tag idea.

\para{{\bf broader implications.}} For production systems, visible thought-token insertion should be treated as a controllable style choice, not a safety or factuality guarantee. Methods with stronger empirical support include retrieval grounding, critique-and-revise pipelines, and explicit external verification.
